\begin{ex}
    (UFRGS) Existem \textit{n} maneiras distintas de marcar 6 quadrados na figura, marcando exatamente dois em cada coluna e um em cada linha. O valor de \textit{n} é:
\begin{center}
\begin{tikzpicture}
 \draw (0,0)--(0,6)--(3,6)--(3,0)--(0,0);
 \draw (1,0)--(1,6); 
 \draw (2,0)--(2,6);
 \draw (0,1)--(3,1);
 \draw (0,2)--(3,2);
 \draw (0,3)--(3,3);
 \draw (0,4)--(3,4);
 \draw (0,5)--(3,5);
\end{tikzpicture}
\end{center}
   \begin{enumerate}[(a)]
   \item 36
   \item 45
   \item 60
   \item 90
   \item 120
   \end{enumerate}
     \begin{sol}
       resposta: d \\
       $\mathrm{C}_{6,2}\cdot\mathrm{C}_{4,2}\cdot\mathrm{C}_{2,2}=90$
     \end{sol}
\end{ex}